\section{Introdução}

Os nematoides fitoparasitas constituem um dos principais fatores limitantes da produtividade agrícola, pois comprometem o desenvolvimento radicular e reduzem a capacidade das plantas de absorver água e nutrientes. Entre as espécies de maior importância para a cultura da soja destaca-se o nematoide de cisto, \textit{Heterodera glycines}, considerado um dos principais agentes de perdas produtivas dessa cultura \cite{bent2022}. A persistência de seus cistos no solo por períodos prolongados dificulta a erradicação do patógeno e reforça a necessidade de métodos capazes de identificar áreas infestadas e apoiar a definição de estratégias de manejo.

O diagnóstico convencional de \textit{H. glycines} baseia-se principalmente na coleta de amostras de solo e raízes, que são encaminhadas a laboratórios especializados para extração de cistos e ovos, contagem e identificação do nematoide. O processamento envolve etapas sucessivas de peneiramento, lavagem, separação e análise microscópica, cuja execução demanda tempo e profissionais capacitados \cite{arjoune2022,mahecha2022,pun2024}. Métodos moleculares e outras abordagens de diagnóstico também podem ser empregados, conforme a finalidade da análise e os recursos disponíveis \cite{shao2023}. Embora os procedimentos laboratoriais sejam essenciais para confirmar a presença e estimar a densidade populacional do nematoide, sua aplicação a grande quantidade de amostras depende de infraestrutura e pessoal capacitado. O intervalo entre coleta e resultado pode variar com o método utilizado e o volume de amostras recebido pelo laboratório; por isso, não se pressupõe um prazo único para todas as situações.

Além da análise laboratorial, o fluxo de diagnóstico envolve coleta, identificação, acondicionamento e transporte das amostras. Em propriedades distantes dos laboratórios especializados, o envio do material pode demandar deslocamentos, transporte contratado ou serviços de entrega e embalagens adequadas. Dependendo das orientações do laboratório, do material coletado e do tempo de deslocamento, podem ser necessários recipientes específicos e cuidados para evitar exposição ao calor ou ressecamento; em alguns protocolos, pode-se recorrer a caixas térmicas. Materiais de acondicionamento, combustível, frete e mão de obra são componentes possíveis do custo operacional, mas sua contribuição depende do contexto local e não foi quantificada neste estudo. Portanto, a eventual economia proporcionada por uma triagem no local deve ser tratada como hipótese para avaliação futura, e não como resultado demonstrado.

O intervalo entre a coleta e o processamento também pode afetar a qualidade das amostras, conforme o material e as condições de armazenamento. Variações de temperatura e umidade, exposição direta à radiação solar, ressecamento ou acondicionamento inadequado são fatores que precisam ser considerados para preservar solo, raízes e organismos presentes. Por essa razão, as amostras devem ser corretamente identificadas, protegidas e encaminhadas segundo as orientações do laboratório responsável. A distância entre a propriedade e o laboratório pode dificultar o controle dessas condições e o acesso rápido ao resultado, mas esse efeito não foi medido no presente experimento. Uma tecnologia capaz de realizar triagem no próprio local de amostragem poderia, se validada em campo, reduzir parte da dependência logística; seu efeito sobre a integridade das amostras transportadas ainda precisaria ser testado.

Outro desafio está relacionado à distribuição espacial do nematoide no campo. Populações de \textit{H. glycines} podem apresentar padrões de agregação e diferenças de densidade dentro de um mesmo talhão \cite{avendano2003,poromarto2019}. Consequentemente, a representatividade do diagnóstico depende da quantidade e da distribuição dos pontos amostrados. O aumento da densidade amostral pode proporcionar uma caracterização mais detalhada da área, mas também amplia a quantidade de recipientes, registros, materiais de acondicionamento, deslocamentos e análises laboratoriais necessários. Assim, o custo operacional tende a aumentar conforme cresce o número de pontos avaliados.

Nesse contexto, uma ferramenta portátil de triagem poderia permitir a realização de medições em um número maior de pontos sem exigir que todas as amostras fossem imediatamente encaminhadas ao laboratório. As leituras realizadas no local poderiam ser utilizadas para identificar regiões com padrões distintos e selecionar os pontos prioritários para análises nematológicas confirmatórias. Essa abordagem não pretende substituir os métodos laboratoriais convencionais, mas complementar o processo de diagnóstico. A possibilidade de reduzir o volume de amostras transportadas e processadas, bem como os custos relacionados à logística, ao acondicionamento e à mão de obra, depende de validação em campo e de uma análise operacional específica, ainda não realizadas neste estudo.

Tecnologias baseadas em sensores têm sido investigadas como alternativas para o monitoramento das condições do solo e das plantas. Sensores de pH, condutividade elétrica, umidade e temperatura, assim como sistemas ópticos de sensoriamento remoto, permitem acompanhar variáveis relacionadas ao desenvolvimento das culturas \cite{santos2022}. Entretanto, essas variáveis apresentam especificidade limitada para a detecção de nematoides, uma vez que diferentes fatores bióticos e abióticos podem produzir respostas semelhantes. A análise de compostos orgânicos voláteis (VOCs) constitui uma abordagem complementar, pois alterações fisiológicas decorrentes da interação entre planta e patógeno podem modificar o perfil volátil do sistema biológico. Perfis de VOCs também respondem a estresses não relacionados ao patógeno investigado; por isso, uma diferença entre grupos experimentais não comprova, isoladamente, especificidade para \textit{H. glycines} \cite{gan2023}.

Alterações na emissão de VOCs de plantas de soja infestadas por \textit{H. glycines} já foram investigadas por cromatografia gasosa acoplada à espectrometria de massas (GC--MS). \citeonline{constantino2021} observaram diferenças entre os perfis voláteis da parte aérea de plantas infestadas e não infestadas, com separação das condições experimentais a partir do quinto dia após a infestação, antes do aparecimento de diferenças significativas no porte das plantas. Esses resultados indicam que a interação entre a soja e o nematoide pode produzir alterações detectáveis na composição volátil. Contudo, o estudo concentrou-se nos compostos emitidos pela parte aérea, mantendo o solo e o sistema radicular isolados durante a coleta. Em outro contexto, \citeonline{velloso2021} mostraram que compostos liberados por diferentes plantas podem modificar a eclosão de ovos de \textit{H. glycines}; esse achado sustenta a relevância biológica dos voláteis na interação planta--nematoide, mas não demonstra um método de diagnóstico. Permanece, portanto, a necessidade de investigar se as condições experimentais com e sem inoculação também podem ser diferenciadas a partir do perfil volátil integrado do sistema solo--planta.

A hipótese deste trabalho é que a inoculação possa alterar, direta ou indiretamente, o \textit{headspace} amostrado no sistema solo--planta. O padrão registrado pode refletir respostas fisiológicas da soja, mudanças na atividade microbiana do substrato, efeitos do material utilizado na inoculação ou alterações gerais do crescimento vegetal. A atividade microbiana do solo, por si, pode gerar sinais detectáveis por narizes eletrônicos \cite{decesare2011}. Assim, a discriminação das classes experimentais não permite identificar a origem química de cada sinal nem concluir que os sensores respondem diretamente ao nematoide.

Embora a GC--MS permita separar e identificar individualmente os compostos presentes em uma amostra, sua aplicação geralmente requer instrumentação especializada, preparação das amostras e operação por pessoal capacitado. O tamanho e as exigências operacionais dos equipamentos também limitam sua utilização em medições rápidas e diretamente em condições de campo \cite{banerji2024}. Em contraste, os sistemas de nariz eletrônico utilizam matrizes de sensores de gases com seletividade parcial, cujas respostas combinadas formam um padrão multivariado associado à composição volátil da amostra. Esse padrão, denominado assinatura olfativa, pode ser processado por técnicas de reconhecimento de padrões e aprendizado de máquina para diferenciar amostras sem a necessidade de identificar individualmente cada composto \cite{berna2010,li2024enose}. Embora não substituam a capacidade analítica da GC--MS, os narizes eletrônicos apresentam potencial para medições rápidas, portáteis e com menor complexidade operacional \cite{li2024enose}. Ao mesmo tempo, apresentam limitações de sensibilidade e especificidade, deriva de sinal e interferência das condições ambientais, razão pela qual devem ser avaliados como ferramentas complementares ao diagnóstico confirmatório \cite{cellini2017}.

Narizes eletrônicos têm sido investigados como ferramentas não invasivas para o reconhecimento de alterações associadas a pragas e doenças em plantas a partir dos VOCs emitidos pelo material vegetal \cite{cui2018,macdougall2022}. A portabilidade e a possibilidade de produzir respostas diretamente no local de amostragem tornam esses sistemas relevantes para aplicações que exigem grande quantidade de medições. Em uma aplicação futura, após validação com infestações naturais e condições ambientais variáveis, a realização de uma triagem em campo poderia reduzir o tempo entre a coleta e a obtenção de uma indicação preliminar, além de diminuir a necessidade de transportar e acondicionar todas as amostras coletadas. O encaminhamento ao laboratório poderia ser direcionado aos pontos considerados prioritários ou aos casos que exigissem confirmação. Tais ganhos operacionais são possibilidades de aplicação e não resultados avaliados neste experimento; a transição de sistemas de sensores para uso em campo permanece um desafio \cite{macdougall2022}.

Entretanto, a aplicação dessa tecnologia à triagem de amostras associadas a \textit{H. glycines} a partir do sistema solo--planta ainda necessita de investigação. Além da condição experimental atribuída aos vasos, fatores como a umidade do solo e as condições empregadas durante a aquisição podem modificar a resposta dos sensores. O teor de água do solo pode alterar a liberação e a retenção de voláteis, enquanto a umidade do gás e a temperatura podem interferir na resposta sensorial \cite{asensio2007,cellini2017}. Neste estudo, a umidade do solo é uma condição experimental e um possível fator de confusão, e sua leitura também pode fornecer informação preditiva ao classificador. A pressão registrada durante a aquisição descreve condições de operação da bomba e de transporte do \textit{headspace} até a câmara sensorial; eventual importância preditiva dessa variável pode refletir o protocolo, não um efeito específico do nematoide. Esses fatores devem ser considerados no desenvolvimento dos protocolos experimentais e dos métodos de classificação, especialmente para que as diferenças identificadas não sejam atribuídas à condição fitossanitária sem examinar a influência das variações ambientais e operacionais.

Diante desse cenário, este estudo avaliou, em ambiente controlado, se os sinais de um nariz eletrônico obtidos durante a coleta a vácuo do \textit{headspace} de vasos de soja permitem discriminar o grupo experimentalmente inoculado com \textit{H. glycines} do grupo controle não inoculado. Também examinou a associação do desempenho com o índice operacional de umidade do solo, a pressão registrada durante a aquisição, a contribuição dos sensores da matriz e o uso de sinais temporais na classificação. Os rótulos representam a alocação experimental dos vasos e não a confirmação nematológica individual ou a densidade populacional do patógeno. Como as coletas foram repetidas nos mesmos vasos e a divisão principal por coleta não separou todos os vasos entre treinamento e teste, os resultados são interpretados como prova de conceito, sem inferência de desempenho diagnóstico em vasos independentes ou em condições de campo. Após validação externa, a abordagem poderá ser avaliada como ferramenta portátil de triagem para ampliar a aquisição de dados no local de amostragem e apoiar o monitoramento espacial de áreas agrícolas. Seu potencial para reduzir custos de coleta, acondicionamento, transporte e processamento laboratorial permanece uma hipótese para estudos posteriores.
